\chapter{Анализ предметной области}
\label{chap:domain-analysis}

\section{Микросервисы и наблюдаемость}

Микросервисная архитектура рассматривает приложение как набор небольших сервисов,
которые выполняются как самостоятельные процессы, взаимодействуют по легковесным
протоколам и могут разворачиваться независимо друг от друга~\cite{fowlerMicroservices,richardsonMicroservicePattern}.
Такой подход упрощает изоляцию бизнес-функций и независимое развитие компонентов,
однако одновременно делает выполнение пользовательского сценария распределенным.
Даже простой HTTP-запрос может пройти через несколько сервисов, каждый из которых
выполняет собственную часть логики, обращается к внешним зависимостям и формирует
свои журналы событий.

Для эксплуатационной диагностики этого уже недостаточно описывать систему только
через отдельные процессы и сетевые вызовы. В распределенной среде требуется видеть
не только факт ошибки в конкретном сервисе, но и причинно-следственную цепочку,
которая к этой ошибке привела. Если один сервис вернул ошибку по таймауту,
оператору важно понимать, был ли источник проблемы во внутренней бизнес-логике,
в нисходящем вызове, в повторной попытке запроса или в несогласованности данных
между несколькими компонентами. Чем больше количество сервисов и точек
интеграции, тем труднее вручную восстановить полный путь запроса по локальным
логам и метрикам.

Наблюдаемость описывает способность понимать внутреннее состояние системы по ее
внешним проявлениям~\cite{otelWhatIs,otelObservabilityPrimer}. В программных
системах это обычно достигается через анализ телеметрических данных: логов,
метрик и трейсов. Чтобы система была наблюдаемой, ее необходимо инструментировать,
то есть встроить в код механизмы генерации телеметрии и передачи ее во внешние
средства анализа~\cite{otelWhatIs}. Для микросервисов наблюдаемость становится не
дополнительной функцией, а обязательным условием эксплуатации, поскольку отказ
одного компонента часто проявляется как деградация совокупного пользовательского
сценария.

Особенность микросервисной архитектуры состоит в том, что многие проблемы
относятся к классу межсервисных. Например, время ответа пользовательского
программного интерфейса может увеличиться не из-за локальной перегрузки, а из-за
каскада обращений к нескольким зависимым сервисам. В такой ситуации необходимо
видеть весь путь обработки запроса: от точки входа в систему до последнего
зависимого вызова. Именно эта потребность делает распределенную трассировку одним
из ключевых элементов наблюдаемости для микросервисов~\cite{otelSignals,jaegerArchitecture}.

\begin{figure}[H]
    \centering
    \begin{tikzpicture}[
        node distance=1.8cm,
        service/.style={draw, rounded corners, minimum width=3.3cm, minimum height=0.9cm, align=center},
        arrow/.style={->, thick}
    ]
        \node[service] (client) {Внешний клиент};
        \node[service, right=2.0cm of client] (gateway) {Входной сервис\\HTTP-интерфейс};
        \node[service, right=2.0cm of gateway, yshift=1.0cm] (serviceb) {Сервис\\обработки заказа};
        \node[service, right=2.0cm of gateway, yshift=-1.0cm] (servicec) {Сервис\\оценки риска};
        \node[service, right=2.0cm of serviceb, yshift=-0.5cm] (storage) {Внешняя система\\или база данных};

        \draw[arrow] (client) -- node[above] {HTTP-запрос} (gateway);
        \draw[arrow] (gateway) -- node[above] {вызов 1} (serviceb);
        \draw[arrow] (gateway) -- node[below] {вызов 2} (servicec);
        \draw[arrow] (serviceb) -- node[above] {доп.\ запрос} (storage);
        \draw[arrow] (servicec) |- (storage);
    \end{tikzpicture}
    \caption{Пример прохождения одного пользовательского запроса через несколько сервисов}
    \label{fig:distributed-request}
\end{figure}

На рисунке~\ref{fig:distributed-request} показан типичный сценарий, в котором один
входящий HTTP-запрос к входному сервису инициирует несколько внутренних вызовов.
С точки зрения пользователя такой сценарий воспринимается как одна операция, но с
технической точки зрения он распадается на набор зависимых действий, выполняемых в
разных процессах. Без единого контекста трассировки сопоставление этих действий
становится трудоемкой ручной задачей.

\section{Логи, метрики и трейсы}

Наблюдаемость в современных системах опирается на несколько взаимодополняющих
типов телеметрии. OpenTelemetry рассматривает их как сигналы, каждый из которых
дает собственный взгляд на поведение системы~\cite{otelSignals,otelOverview}. На
практике для эксплуатационного анализа чаще всего используются логи, метрики и
трейсы.

Логи представляют собой событийные записи о фактах, произошедших в приложении.
Они удобны для фиксации ошибок, бизнес-событий, параметров запросов и произвольных
контекстных полей. Сильная сторона логов состоит в высокой выразительности:
разработчик может записать практически любые данные, которые считает полезными.
Слабая сторона состоит в том, что локальные журналы плохо восстанавливают
межсервисную причинность, если в них не добавлены согласованные идентификаторы
запросов или контекст трассировки.

Метрики описывают систему в виде числовых измерений и агрегатов. Они позволяют
эффективно отвечать на вопросы о частоте ошибок, пропускной способности, времени
ответа, загрузке ресурсов и других измеряемых характеристиках~\cite{otelSignals}.
Метрики особенно полезны для мониторинга и построения оповещений. Однако они
агрегированы по своей природе и потому слабо подходят для анализа конкретного
экземпляра пользовательского запроса.

Трейсы предназначены для представления пути выполнения отдельной операции через
набор сервисов и внутренних шагов~\cite{otelWhatIs,jaegerArchitecture}. Они дают
временную и причинную структуру исполнения: позволяют увидеть, какие вызовы были
выполнены, в каком порядке, сколько времени занял каждый шаг и где возникла
задержка или ошибка. Для распределенных HTTP-сценариев именно трейсы дают ответ на
вопрос, почему операция пользователя завершилась медленно или некорректно.

\begin{table}[H]
    \centering
    \caption{Сравнение логов, метрик и трейсов}
    \label{tab:signals-comparison}
    \renewcommand{\arraystretch}{1.2}
    \begin{tabularx}{\textwidth}{>{\raggedright\arraybackslash}p{2.4cm} >{\raggedright\arraybackslash}X >{\raggedright\arraybackslash}X >{\raggedright\arraybackslash}X}
        \toprule
        Вид данных & Основное назначение & Преимущества & Ограничения \\
        \midrule
        Логи & Фиксация событий, ошибок, бизнес-действий и произвольных полей & Высокая детализация; гибкость формата; удобны для расследования отдельных событий & Плохо показывают полную цепочку распределенного запроса без дополнительной корреляции \\
        \addlinespace
        Метрики & Количественная оценка состояния системы и построение мониторинга & Эффективны для оповещений, анализа тенденций и оценки емкости & Недостаточно контекста для разбора конкретного запроса или причины ошибки \\
        \addlinespace
        Трейсы & Восстановление пути выполнения одной операции через несколько компонентов & Показывают причинность, длительности и структуру вызовов; удобны для анализа задержек & Требуют корректной инструментировки и передачи контекста между сервисами \\
        \bottomrule
    \end{tabularx}
\end{table}

Таблица~\ref{tab:signals-comparison} показывает, что перечисленные сигналы не
конкурируют друг с другом, а решают разные аналитические задачи. Для
микросервисных приложений особенно важна совместная работа всех трех видов данных.
Метрики помогают понять, что в системе началась деградация, логи уточняют
конкретные обстоятельства ошибки, а трейсы показывают, в каком именно участке
распределенного сценария возникла проблема.

\section{Основные понятия распределенной трассировки}

Ключевым объектом распределенной трассировки является трейс. Трейс представляет
собой логическую запись о выполнении одной распределенной операции от точки входа
до завершения всех связанных шагов~\cite{otelOverview,otelSdkTrace}. Обычно трейс
соответствует обработке одного пользовательского запроса, сообщения или фоновой
задачи.

Трейс состоит из спанов. Спан описывает отдельную операцию внутри общего пути
выполнения: прием входящего запроса, обращение к базе данных, вызов внешнего HTTP
сервиса, сериализацию ответа или другой значимый шаг~\cite{otelSdkSpan}. Каждый
спан имеет начало и конец во времени, имя операции и набор метаданных. Спаны
могут быть вложенными, образуя дерево или, в более общем случае, ориентированный
граф причинных связей.

Для идентификации трейса и его элементов используются два ключевых значения:
идентификатор трейса и идентификатор спана. Идентификатор трейса является общим
для всех спанов, принадлежащих одному трейсу, поэтому он позволяет связывать
события из разных сервисов в одну цепочку. Идентификатор спана определяет
конкретный спан внутри этого трейса. Дополнительно спан обычно содержит ссылку на
родительский спан, что позволяет восстанавливать иерархию вызовов~\cite{w3cTraceContext,otelTraceApi}.

Практически важны также атрибуты и события спана. Атрибуты представляют собой
структурированные пары ключ--значение, которые описывают свойства операции:
HTTP-метод, маршрут, код ответа, имя сервиса, адрес удаленного узла и другие
характеристики. События спана фиксируют значимые моменты внутри операции, которые
произошли между ее началом и завершением. За счет этого трейс дает не только
временную диаграмму вызовов, но и сопровождающий контекст, необходимый для
диагностики~\cite{otelOverview,otelTraceApi}.

Важное свойство распределенной трассировки состоит в сохранении связи между
входящими и исходящими операциями. Если сервис принимает HTTP-запрос, а затем
вызывает другой сервис, то исходящий клиентский спан должен быть логически связан
с входящим серверным спаном. При корректной инструментировке это позволяет видеть
всю цепочку вызовов как единый трейс, а не как набор независимых локальных
операций.

\section{Стандарт W3C для контекста трассировки}

Для межсервисной трассировки недостаточно локально создавать спаны. Необходимо,
чтобы разные процессы одинаково понимали, как передавать контекст текущего трейса в
сетевых сообщениях. Эту задачу решает стандарт W3C Trace Context~\cite{w3cTraceContext}.
Он определяет заголовки HTTP и формат значений, которые позволяют переносить
контекст трассировки между сервисами, прокси и другими участниками распределенного
взаимодействия.

Основным заголовком является `traceparent`. Он содержит версию формата,
идентификатор трейса, идентификатор текущего родительского спана и флаги
трассировки~\cite{w3cTraceContext}. Благодаря этому вызываемый сервис может не
начинать новый независимый трейс, а продолжить уже существующий. Второй
заголовок, `tracestate`, предназначен для передачи дополнительной информации,
специфичной для конкретных поставщиков телеметрической инфраструктуры или
механизмов обработки данных трассировки~\cite{w3cTraceContext}.

Стандартизация формата контекста имеет принципиальное значение. В распределенной
системе один сервис может быть написан на Rust, другой на Go, третий на Java, а
сетевые вызовы могут проходить через промежуточные шлюзы и балансировщики.
Если каждый компонент использует собственный формат идентификаторов и заголовков,
единый трейс распадается. Использование W3C Trace Context позволяет обеспечить
совместимость между библиотеками, каркасами приложений и системами
наблюдаемости.

Для HTTP-сервисов поддержка `traceparent` особенно важна, поскольку именно этот
заголовок обычно становится минимальным переносимым контекстом между входящим и
исходящим запросом. Следовательно, библиотека распределенной трассировки для Rust
должна не только создавать спаны внутри процесса, но и корректно извлекать и
внедрять контекст трассировки W3C в HTTP-заголовки.

\section{OpenTelemetry, OTLP и Jaeger}

OpenTelemetry представляет собой открытый вендор-нейтральный стандарт и набор
инструментов для генерации, сбора и экспорта телеметрических данных, включая
трейсы, метрики и журналы событий~\cite{otelDocs,otelWhatIs}. Для темы настоящей
работы важны три роли, которые OpenTelemetry разделяет на уровне архитектуры.

Первая роль --- инструментирование приложений. На этом уровне код приложения или
библиотека создают спаны, события, атрибуты и контекст распространения. Вторая
роль --- передача и экспорт телеметрии. Для этого в экосистеме OpenTelemetry
используется протокол OTLP, который описывает формат и механизмы доставки
телеметрических данных между источниками, промежуточными узлами и системами
хранения и визуализации~\cite{otlpSpec,otlpExporterSpec}. Третья роль --- система
наблюдаемости, в которой данные сохраняются, агрегируются и визуализируются.

OTLP является базовым протоколом OpenTelemetry для переноса телеметрических
данных. Спецификация определяет кодирование, транспорт и общие требования к
доставке сигналов~\cite{otlpSpec}. Для данных трассировки на практике часто
используется OTLP/HTTP, поскольку он хорошо сочетается с типичным стеком
HTTP-приложений и поддерживается многими комплектами средств разработки и
сборщиками телеметрии~\cite{otlpExporterSpec}. В рассматриваемом классе задач
OTLP/HTTP выступает как транспорт между приложением и сборщиком телеметрии.

Jaeger является специализированной системой распределенной трассировки, которая
обеспечивает прием, хранение и визуализацию данных трассировки~\cite{jaegerDocs,jaegerArchitecture}.
С точки зрения архитектуры наблюдаемости Jaeger удобно рассматривать как систему
для просмотра и анализа трейсов. Пользователь может открыть трейс, увидеть дерево
спанов, длительности операций, временные диаграммы и атрибуты отдельных шагов.
Это делает Jaeger подходящим инструментом для проверки корректности передачи
контекста и для поиска узких мест в распределенных HTTP-сценариях.

Во многих практических схемах приложение не отправляет данные непосредственно в
Jaeger, а использует сборщик OpenTelemetry как промежуточный узел. Сборщик
телеметрии принимает OTLP-данные от приложений, может выполнять маршрутизацию,
преобразование или буферизацию и затем передает данные трассировки в конечную
систему хранения и визуализации~\cite{jaegerArchitecture,otlpSpec}. Такая схема
уменьшает связанность приложения с конкретной системой визуализации и упрощает
развитие инфраструктуры наблюдаемости.

\section{Экосистема Rust для трассировки HTTP-сервисов}

В Rust уже существует набор зрелых компонентов, необходимых для построения
распределенной трассировки, однако они покрывают разные уровни задачи и потому
требуют согласованной интеграции.

Библиотека `tracing` предоставляет базовые интерфейсы для инструментирования
приложений и библиотек, то есть для создания спанов, событий и структурированных
полей диагностики~\cite{tracingDocs,tracingSubscriberDocs}. Она решает задачу
локального описания операций внутри процесса, но сама по себе не задает единый
стандарт контекста трассировки и не определяет систему экспорта.

Библиотека `opentelemetry` и связанный с ней набор средств разработки для Rust дают
интерфейсы и реализации для работы с контекстом, средствами передачи контекста,
поставщиком трассировщиков и другими сущностями экосистемы OpenTelemetry~\cite{opentelemetryPropagationDocs,opentelemetrySdkTraceDocs}.
Через эти компоненты реализуются правила распространения контекста и экспорта
данных трассировки во внешнюю систему наблюдаемости.

Библиотека `tower` задает фундаментальные абстракции `Service` и `Layer`, которые
используются при построении прослоек в асинхронных сетевых приложениях~\cite{towerDocs,towerServiceDocs,towerLayerDocs}.
Эти абстракции удобны для внедрения серверной трассировки на границе обработки
HTTP-запроса, поскольку прослойка может перехватить входящее сообщение, извлечь
контекст, создать спан и передать управление внутреннему обработчику.

Каркас `axum` предоставляет маршрутизацию и обработку HTTP-запросов, опираясь на
экосистему Tower~\cite{axumDocs}. Благодаря этому для Axum естественно строить
решение в виде слоя промежуточной обработки, который добавляется к маршрутизатору
и автоматически инструментирует входящие запросы.

\begin{table}[H]
    \centering
    \caption{Роль основных Rust-компонентов в задаче распределенной трассировки}
    \label{tab:rust-components}
    \renewcommand{\arraystretch}{1.2}
    \begin{tabularx}{\textwidth}{>{\raggedright\arraybackslash}p{2.7cm} >{\raggedright\arraybackslash}X >{\raggedright\arraybackslash}X}
        \toprule
        Компонент & Основная роль & Что не закрывает самостоятельно \\
        \midrule
        `tracing` & Создание спанов, событий и структурированных диагностических записей & Не задает полный конвейер экспорта и межсервисную передачу контекста \\
        \addlinespace
        `opentelemetry` и `opentelemetry\_sdk` & Контекст, средства передачи контекста, поставщик трассировщиков, интеграция со стандартом OpenTelemetry & Требуют явной настройки и встраивания в серверный и клиентский HTTP-код \\
        \addlinespace
        `tower` & Абстракции `Layer` и `Service` для промежуточных слоев & Не реализует логику наблюдаемости сам по себе \\
        \addlinespace
        `axum` & HTTP-каркас для построения сервисов на основе Tower & Не дает готовой интеграции распределенной трассировки из коробки \\
        \bottomrule
    \end{tabularx}
\end{table}

Таблица~\ref{tab:rust-components} показывает, что перечисленные компоненты хорошо
подходят для построения решения распределенной трассировки, но относятся к разным
уровням абстракции. Разработчик должен самостоятельно связать локальную
инструментацию, прослойку для входящих запросов, передачу контекста для исходящих
HTTP-вызовов, экспорт данных трассировки по OTLP и корреляцию журналов событий с
идентификатором трейса.

\section{Проблема ручной интеграции}

Несмотря на зрелость перечисленных стандартов и библиотек, их наличие не означает
автоматического получения простого подключения распределенной трассировки в
сервисе на Rust. На практике разработчику приходится вручную решать несколько
задач одновременно.

Во-первых, необходимо инициализировать общий конвейер трассировки приложения:
настроить подписчик, поставщик трассировщиков, средство передачи контекста и
экспортер. Во-вторых, требуется инструментировать входящие HTTP-запросы на
границе сервиса, чтобы корректно извлекать входящий `traceparent`, создавать
серверный спан и фиксировать атрибуты запроса и ответа. В-третьих, нужно отдельно
обрабатывать исходящие HTTP-вызовы, внедряя контекст трассировки в заголовки
запросов к вызываемым сервисам и создавая клиентский спан. Наконец, полезно
коррелировать журналы стандартного вывода с текущим трейсом, чтобы события
журналов можно было сопоставлять с визуализацией в системе просмотра трейсов.

По отдельности каждая из этих задач решаема стандартными средствами
экосистемы Rust, однако совокупная интеграция остается трудоемкой и подверженной
ошибкам. Проблемы возникают не только на уровне кода, но и на уровне
согласованности: необходимо, чтобы одни и те же правила именования, передачи
контекста и экспорта применялись во всех сервисах проекта. При отсутствии единого
библиотечного слоя разные команды или разные сервисы легко приходят к
несовместимым схемам настройки наблюдаемости.

Следовательно, для микросервисов на Rust существует практическая потребность в
библиотеке, которая предоставляет единый простой путь подключения распределенной
трассировки. Такая библиотека должна скрывать повторяющуюся инфраструктурную
настройку, автоматически поддерживать стандартный контекст трассировки W3C и интегрироваться со
стандартными для Rust HTTP-компонентами. Этот вывод задает основание для
последующих разделов работы, в которых рассматриваются модель, архитектура и
реализация специализированной библиотеки распределенной трассировки.
